\color{unimain}{\subsection{Cигнализация}}
\color{black}

\begin{enumerate}

\item
Устройство формирует выходное воздействие на сигнализацию «Срабатывание» по схеме «ИЛИ» в следующих случаях:
\begin{itemize}
\item Срабатывание ДЗТ;
\item Срабатывание газовых и технологических защит;
\item Срабатывание МТЗ НН;
\item Срабатывание ЗПО;
\item Срабатывание ЗП;
\item Срабатывание УРОВ;
\item Срабатывание внешних защит.
\end{itemize}

\item
Устройство формирует выходное воздействие на сигнализацию «Неисправность» по схеме «ИЛИ» в следующих случаях:
\begin{itemize}
\item Неисправность цепей тока ДЗТ;
\item Снижение изоляции цепей газовых и технологических защит;
\item Минимальный/максимальный уровень масла защищаемого оборудования;
\item Отказ/неисправность системы охлаждения АТ;
\item Неисправность цепей напряжения стороны НН;
\item Неисправность силового выключателя;
\item Снижение изоляции ошиновки НН;
\item Нарушение режима по испытательным блокам и/или переключателям в выходных цепях;
\item Потеря GOOSE-сообщения;
\item Срабатывание внешней функции на сигнализацию;
\item Неисправность цепей оперативного тока.
\end{itemize}

\item
Алгоритм работы логики логики предупредительной сигнализации приведен на рисунке \ref{pic:SIGN}.

\begin{figure}[H]
\centering
\includegraphics[width=1\textwidth,height=1\textheight,keepaspectratio]{\fbpath/06.02 ДЗАТ/91 Сигнализация/img/SIGN.pdf}
\captionof{figure}{Функциональная схема логики блока сигнализации}\label{pic:SIGN}
\end{figure}

\end{enumerate}